\documentclass{article}
\usepackage[utf8]{inputenc}
\usepackage[margin=0.8in]{geometry}
\usepackage{graphicx}
\usepackage{pgffor}
\usepackage{caption}
\usepackage{paracol}
\usepackage{chngcntr}
\usepackage{amsmath}
\usepackage{biblatex}
%\counterwithout{figure}{section} % If using Article

\addbibresource{references.bib}


\newcommand{\singleWidth}{0.6\textwidth}

\captionsetup[figure]{
    font=small,           % Makes text smaller (can use 'footnotesize' for even smaller)
    labelfont=bf,         % Makes "Figure X" bold
    margin=10pt,          % Indents caption from the left/right of the column
    skip=5pt,             % Reduces space between image and caption
    justification=centering % Keeps it neat
}


\title{Final Project \& Literature Review Proposal}
\author{Group 17: \\ Darlington Nkrumah (202492437) \\
  Greg de Souza (202582225) \\
  Xuan Toan Doan (202583882)}
\date{February 2026}

\begin{document}

\maketitle


\section*{Final Project Proposal: Terrain classification
  for NASA's mars2020 (Perseverance) rover}

\textbf{AI4Mars} is a collection of open dataset of images taken by the various
rovers currently operating on mars. The images are also label by
citizen science initiatives through the \textit{zooniverse} platform,
with the final labels being verified by experts. The dataset we chose to
work with is the navigation images of the Mars 2020 Perseverance rover.

The dataset consists of 3827 unique labeled images taken by the navigation cameras with different
portions of the terrain labeled into one of 5 classes, according to table
\ref{tab:labels}. The labels are written into the rgb colors of a label image.
A reference of representative classes is show in figure \ref{Fig:terrain},
and an example of the label mask is shown in figure \ref{Fig:label}. This images
are taken from a conference paper about the AI4Mars project \cite{AI4Mars_2021} and
are from the MSL Curiosity navigational dataset. Our dataset is more recent and from
a different rover, but the classes and classification methods are the same.


\begin{table}[]
  \centering
  \begin{tabular}{l|l|l}
    \textbf{Label} & \textbf{Type}             & \textbf{RGB Label}         \\
    0     & soil             & {[}0,0,0{]}       \\
    1     & bedrock          & {[}1,1,1{]}       \\
    2     & sand             & {[}2,2,2{]}       \\
    3     & big rock         & {[}3,3,3{]}       \\
    255   & null / unlabeled & {[}255,255,255{]}
  \end{tabular}
  \caption{Labels for the Mars2020 Navigation dataset}
  \label{tab:labels}
\end{table}

\begin{figure}[htbp]
  \centering
  % Includes the image, scaled to half the text width
  \includegraphics[width=0.5\textwidth]{figures/TerrainTypes.png} 
  \caption{Representative examples of each class, taken from \cite{AI4Mars_2021}}
  \label{Fig:terrain}
\end{figure}

\begin{figure}[htbp]
  \centering
  \includegraphics[width=0.8\textwidth]{figures/LabeledImages.png} 
  \caption{Example of original navigational camera image vs labeled image,
    taken from \cite{AI4Mars_2021}}}
  \label{Fig:label}
\end{figure}

The goal of our final project is to learn how to use the images provided by the
AI4Mars data set to generate a design matrix, apply one or more machine learning
method to segment the images in the 4 classes, and compare our results with other
published results for martian terrain classification \cite{AI4Mars_2022}
\cite{AI4Mars_2025}.


\section*{Literature Review}

Our literature review will entail studying two different sets of articles.
The first are the articles related to the AI4mars project, detailed the dataset
labeling process, as well as given an overview of the available data, and then assess
the ability of different Machine Learning Models to do image segmentation on the images
provided in the AI4Mars dataset or similar datasets\cite{AI4Mars_2021}\cite{AI4Mars_2022}\cite{AI4Mars_2025}\cite{Atha_Swan_Didier_Hasnain_Ono_2022}\cite{Goutham_Juneja_C_Prasad_2022}\cite{Zhang_Lin_Fan_Wang_Liu_2024}.

The other set of papers we will study will be related to other image segmentation methods unrelated
to the AI4Mars project, with the goal to inform which data processing techniques and methods we
could apply in our final project
\cite{Ghosh_Das_Das_Maulik_2019}\cite{Kale_Thorat_2021}\cite{Minaee_Boykov_Porikli_Plaza_Kehtarnavaz_Terzopoulos_2020}\cite{Sultana_Sufian_Dutta_2020}.



\printbibliography

\end{document}
